\chapter{Foundation I: Efficient Inference for Noisy LLM-as-a-Judge Evaluation} \label{chap:foundation_llm}

This chapter summarizes the first methodological foundation used in the thesis, the preprint \emph{Efficient Inference for Noisy LLM-as-a-Judge Evaluation} \parencite{llmjudge_eif}. Although that paper is motivated by LLM-as-a-judge evaluation, its core statistical structure is directly relevant here: a large set of cheap surrogate labels is available for many units, while a smaller calibration set contains both the surrogate and the gold-standard outcome.

\section{Surrogate Labels and Calibration Design}

Let $Y$ denote the latent target outcome and let $\widehat{Y}$ denote a noisy surrogate or judge label. The data are split into two parts:
\begin{enumerate}
    \item a large evaluation set with only $\widehat{Y}$ observed, and
    \item a smaller calibration set with both $Y$ and $\widehat{Y}$ observed.
\end{enumerate}
The inferential goal is to estimate a mean-type target parameter,
\[
\theta = \mathbb{E}[Y],
\]
while using the large surrogate-only sample to reduce variance and the small calibration sample to correct bias.

The paper compares two broad correction strategies. The first treats $\widehat{Y}$ as a mismeasured version of $Y$ and uses direct measurement-error correction, such as Rogan--Gladen-style estimators. The second treats $\widehat{Y}$ as a generic surrogate outcome and uses prediction-powered correction based on calibration residuals. The central message is that these views can be unified through efficient influence functions and semiparametric efficiency \parencite{rogan_gladen, eif_ref, llmjudge_eif}.

\section{Relationship to Earlier Statistical Literature}

The conceptual roots of this chapter lie in two older literatures. The first is the literature on imperfect diagnostics and measurement error, where the observed label is a noisy proxy for an unobserved truth. Rogan--Gladen-style prevalence correction is a classical example of this perspective \parencite{rogan_gladen}. The second is semiparametric efficiency theory, where efficient influence functions characterize the best possible first-order use of auxiliary information under a given model \parencite{eif_ref}. The LLM-as-a-judge paper is important because it connects these traditions to modern AI-assisted evaluation problems rather than treating them as separate worlds.

That connection matters for the rest of the thesis. In protein-edge prediction, the cheap label is not a binary judge produced by an LLM but a model-based prediction score for an interaction edge. Even so, the statistical skeleton is the same: the inexpensive signal is observed broadly, the expensive signal is observed only on a smaller labeled subset, and the inferential problem is how to combine them efficiently.

\section{Why This Paper Matters for the Thesis}

The importance of \textcite{llmjudge_eif} for this thesis is conceptual rather than domain-specific. In the protein application, the surrogate is not an LLM judge score but a learned edge-prediction score produced by a neural model. The structure is nevertheless the same:
\begin{enumerate}
    \item the surrogate is available at scale,
    \item the gold label is expensive, and
    \item the calibration set is used to correct systematic error.
\end{enumerate}

The paper also clarifies a useful modeling distinction. A cheap surrogate can be treated as a noisy measurement, but it can also be treated as a generic prediction that is corrected through calibration. The latter perspective is more flexible for this thesis because the protein-edge surrogate is continuous, model-based, and not naturally parameterized by a simple sensitivity/specificity pair. That flexibility is one reason prediction-powered methods are a better fit for the protein application than a strict diagnostic-test abstraction.

\section{Efficient Estimation Viewpoint}

Prediction-powered inference can be written as a residual-correction estimator. In the simplest mean-setting,
\[
\widehat{\theta}_{\mathrm{PPI}} = \overline{f}_{U} + \frac{1}{n} \sum_{i \in L} \{Y_i - f_i\},
\]
where $f_i$ is the surrogate prediction, $L$ denotes the labeled calibration set, and $U$ denotes the large unlabeled set. PPI++ introduces a tuning parameter $\lambda$,
\[
\widehat{\theta}_{\lambda} = \overline{Y}_{L} + \lambda\, (\overline{f}_{U} - \overline{f}_{L}),
\]
which can improve efficiency when the surrogate is strong \parencite{ppi_framework, ppi_pp}.

The contribution of \textcite{llmjudge_eif} is to place these estimators in an efficient-influence-function framework. That perspective explains when simple correction rules are efficient, when they are suboptimal, and why a stronger surrogate typically leads to shorter confidence intervals after calibration. For the purposes of this thesis, the main takeaway is that better prediction quality helps twice: it improves point prediction, and it also improves downstream inferential efficiency.

\subsection{Key Formulas for Mean Estimation with a Surrogate Label}

The efficient-influence-function perspective can be written explicitly. Let $R$ be the indicator that the gold-standard label is observed, with $\Pr(R=1)=\pi$, and suppose $R$ is independent of $(Y,\widehat{Y})$. Define
\[
\mu(\widehat{Y}) = \mathbb{E}[Y \mid \widehat{Y}].
\]
Then the efficient influence function for the mean target $\theta=\mathbb{E}[Y]$ takes the form
\[
\phi(O)
=
\mu(\widehat{Y}) - \theta
+
\frac{R}{\pi}\{Y-\mu(\widehat{Y})\},
\]
where $O=(R,RY,\widehat{Y})$. The corresponding efficient estimator therefore has the generic residual-correction form
\[
\widehat{\theta}(\mu)
=
\frac{1}{N}\sum_{i \in U}\mu(\widehat{Y}_i)
+
\frac{1}{n}\sum_{i \in L}\{Y_i-\mu(\widehat{Y}_i)\}.
\]
This equation is the bridge between the semiparametric theory and the practical estimators used later in the thesis. Standard PPI is obtained by the working choice $\mu(\widehat{Y}) \approx \widehat{Y}$, whereas PPI++ uses the linear approximation $\mu(\widehat{Y}) \approx \lambda \widehat{Y}$ \parencite{ppi_framework, ppi_pp, llmjudge_eif}. Detailed algebra for these identities is summarized in Appendix~\ref{chap:appendix-a}.

\subsection{Binary Misclassification View and Rogan--Gladen}

The direct measurement-error view can also be written explicitly in the binary case. Let
\[
q_0 = \Pr(\widehat{Y}=0 \mid Y=0),
\qquad
q_1 = \Pr(\widehat{Y}=1 \mid Y=1),
\qquad
p = \Pr(\widehat{Y}=1).
\]
Then
\[
p
=
\Pr(\widehat{Y}=1 \mid Y=0)\Pr(Y=0)
+
\Pr(\widehat{Y}=1 \mid Y=1)\Pr(Y=1)
=
(1-q_0)(1-\theta)+q_1\theta,
\]
which simplifies to
\[
p = (q_0+q_1-1)\theta + (1-q_0).
\]
Solving for $\theta$ yields the Rogan--Gladen identity
\[
\theta = \frac{p+q_0-1}{q_0+q_1-1}.
\]
This is the classical correction formula against which residual-correction estimators such as PPI and PPI++ are compared in the LLM-as-a-judge paper \parencite{rogan_gladen, llmjudge_eif}. In the protein chapters, I do not use this binary misclassification model directly, but it remains a useful conceptual reference point.

\section{Bridge to the Protein Application}

The rest of the thesis uses the same general template in a different scientific setting. The protein-edge predictor produces the surrogate, the experimentally labeled subset provides calibration, and prediction-powered procedures deliver corrected summaries with uncertainty quantification. The next chapter adds the second methodological foundation needed for that pipeline: labeled-sample planning and power analysis under prediction-powered inference \parencite{ppi_power}.
